%%%%%%%%%%%%%%%%%%%%%%%%%%%%%%%%%%%%%%%%%%%%%%%%%%%%%%%%%%%%
% 6.6 EXTREMAL COMBINATORICS
% The Composition of Advanced Mathematics
%%%%%%%%%%%%%%%%%%%%%%%%%%%%%%%%%%%%%%%%%%%%%%%%%%%%%%%%%%%%

\section{Extremal Combinatorics}
\label{sec:extremal-combinatorics}

\prerequisite{
Basic counting, set systems, graph theory terminology, binomial
coefficients, partially ordered sets, probability, and familiarity with
combinatorial classes.
}

\learningobjectives{
By the end of this section, the reader should be able to:
\begin{itemize}
    \item explain the central questions of extremal combinatorics;
    \item formulate extremal problems using forbidden configurations;
    \item work with extremal functions;
    \item understand the structure of triangle-free graphs;
    \item state and apply Mantel's theorem;
    \item state and interpret Tur\'an's theorem;
    \item understand Tur\'an graphs;
    \item recognize the role of complete multipartite graphs in extremal
          graph theory;
    \item understand the Erd\H{o}s--Stone theorem as an asymptotic
          generalization;
    \item analyze intersecting families of sets;
    \item state and apply the Erd\H{o}s--Ko--Rado theorem;
    \item work with chains and antichains in the Boolean lattice;
    \item state and apply Sperner's theorem;
    \item understand the Lubell--Yamamoto--Meshalkin inequality;
    \item recognize compression and shifting methods;
    \item use averaging and double counting in extremal arguments;
    \item understand the probabilistic method as an existence tool;
    \item distinguish exact extremal results from asymptotic results;
    \item recognize stability phenomena in extremal combinatorics.
\end{itemize}
}

%%%%%%%%%%%%%%%%%%%%%%%%%%%%%%%%%%%%%%%%%%%%%%%%%%%%%%%%%%%%
\subsection{What Is Extremal Combinatorics?}
\label{subsec:what-is-extremal-combinatorics}
%%%%%%%%%%%%%%%%%%%%%%%%%%%%%%%%%%%%%%%%%%%%%%%%%%%%%%%%%%%%

Enumerative combinatorics asks

\[
\boxed{
\text{How many structures are there?}
}
\]

Extremal combinatorics asks a different kind of question:

\[
\boxed{
\text{How large or small can a structure be under given restrictions?}
}
\]

Typical problems have the form

\[
\boxed{
\max\{|X|:X\text{ satisfies property }P\}
}
\]

or

\[
\boxed{
\min\{|X|:X\text{ satisfies property }P\}.
}
\]

The restrictions often forbid certain substructures.

Thus extremal combinatorics studies the interaction between

\[
\boxed{
\text{size}
+
\text{structure}
+
\text{forbidden configurations}.
}
\]

%%%%%%%%%%%%%%%%%%%%%%%%%%%%%%%%%%%%%%%%%%%%%%%%%%%%%%%%%%%%
\subsection{Extremal Functions}
\label{subsec:extremal-functions}
%%%%%%%%%%%%%%%%%%%%%%%%%%%%%%%%%%%%%%%%%%%%%%%%%%%%%%%%%%%%

An extremal function records the maximum size of an object avoiding a
specified configuration.

In graph theory, if \(H\) is a fixed graph, define

\[
\boxed{
\operatorname{ex}(n,H)
}
\]

to be the maximum number of edges in an \(n\)-vertex graph containing no
copy of \(H\).

The notation

\[
\operatorname{ex}(n,H)
\]

is read ``the extremal number of \(H\) on \(n\) vertices.''

Thus,

\[
\boxed{
\operatorname{ex}(n,H)
=
\max\{
|E(G)|:
|V(G)|=n,\,
H\not\subseteq G
\}.
}
\]

%%%%%%%%%%%%%%%%%%%%%%%%%%%%%%%%%%%%%%%%%%%%%%%%%%%%%%%%%%%%
\subsection{Forbidden Substructures}
\label{subsec:forbidden-substructures-extremal}
%%%%%%%%%%%%%%%%%%%%%%%%%%%%%%%%%%%%%%%%%%%%%%%%%%%%%%%%%%%%

A common extremal problem begins by specifying a forbidden object.

For example:

\begin{itemize}
    \item graphs with no triangle;
    \item set families with no containment;
    \item families in which every two sets intersect;
    \item hypergraphs avoiding a fixed configuration;
    \item sequences avoiding a specified pattern.
\end{itemize}

The central question is then:

\[
\boxed{
\text{How large can the structure be without creating the forbidden
configuration?}
}
\]

%%%%%%%%%%%%%%%%%%%%%%%%%%%%%%%%%%%%%%%%%%%%%%%%%%%%%%%%%%%%
\subsection{Triangle-Free Graphs}
\label{subsec:triangle-free-graphs}
%%%%%%%%%%%%%%%%%%%%%%%%%%%%%%%%%%%%%%%%%%%%%%%%%%%%%%%%%%%%

A triangle in a graph is a copy of

\[
K_3.
\]

A graph is triangle-free if it contains no \(K_3\).

The simplest extremal graph problem asks:

\[
\boxed{
\text{How many edges can an }n\text{-vertex triangle-free graph have?}
}
\]

The answer is given by Mantel's theorem.

%%%%%%%%%%%%%%%%%%%%%%%%%%%%%%%%%%%%%%%%%%%%%%%%%%%%%%%%%%%%
\subsection{Mantel's Theorem}
\label{subsec:mantels-theorem}
%%%%%%%%%%%%%%%%%%%%%%%%%%%%%%%%%%%%%%%%%%%%%%%%%%%%%%%%%%%%

\begin{theorem}[Mantel's Theorem]
Every triangle-free graph on \(n\) vertices has at most

\[
\boxed{
\left\lfloor\frac{n^2}{4}\right\rfloor
}
\]

edges.

Equivalently,

\[
\boxed{
\operatorname{ex}(n,K_3)
=
\left\lfloor\frac{n^2}{4}\right\rfloor.
}
\]
\end{theorem}

The bound is attained by a complete bipartite graph with the two parts
as equal in size as possible.

%%%%%%%%%%%%%%%%%%%%%%%%%%%%%%%%%%%%%%%%%%%%%%%%%%%%%%%%%%%%
\subsection{Balanced Complete Bipartite Graph}
\label{subsec:balanced-complete-bipartite}
%%%%%%%%%%%%%%%%%%%%%%%%%%%%%%%%%%%%%%%%%%%%%%%%%%%%%%%%%%%%

Let the vertex set be divided into two parts of sizes

\[
\left\lfloor\frac n2\right\rfloor
\]

and

\[
\left\lceil\frac n2\right\rceil.
\]

Connect every vertex in one part to every vertex in the other.

The resulting graph is

\[
\boxed{
K_{\lfloor n/2\rfloor,\lceil n/2\rceil}.
}
\]

It is triangle-free because a bipartite graph contains no odd cycle.

Its number of edges is

\[
\left\lfloor\frac n2\right\rfloor
\left\lceil\frac n2\right\rceil
=
\boxed{
\left\lfloor\frac{n^2}{4}\right\rfloor.
}
\]

Thus the bound in Mantel's theorem is sharp.

%%%%%%%%%%%%%%%%%%%%%%%%%%%%%%%%%%%%%%%%%%%%%%%%%%%%%%%%%%%%
\subsection{Proof of Mantel's Theorem by Degrees}
\label{subsec:mantel-proof-degrees}
%%%%%%%%%%%%%%%%%%%%%%%%%%%%%%%%%%%%%%%%%%%%%%%%%%%%%%%%%%%%

Let

\[
G=(V,E)
\]

be triangle-free.

For every edge

\[
uv\in E,
\]

the neighborhoods of \(u\) and \(v\) cannot share another vertex,
otherwise a triangle would occur.

Therefore,

\[
\deg(u)+\deg(v)\leq n.
\]

Sum this inequality over every edge:

\[
\sum_{uv\in E}
\left(
\deg(u)+\deg(v)
\right)
\leq
n|E|.
\]

The left side equals

\[
\sum_{v\in V}\deg(v)^2.
\]

Hence,

\[
\sum_{v\in V}\deg(v)^2
\leq
n|E|.
\]

By the Cauchy--Schwarz inequality,

\[
\left(
\sum_{v\in V}\deg(v)
\right)^2
\leq
n
\sum_{v\in V}\deg(v)^2.
\]

Since

\[
\sum_{v\in V}\deg(v)=2|E|,
\]

we obtain

\[
4|E|^2
\leq
n^2|E|.
\]

Assuming \(|E|>0\), divide by \(|E|\):

\[
4|E|\leq n^2.
\]

Therefore,

\[
\boxed{
|E|\leq\frac{n^2}{4}.
}
\]

Since \(|E|\) is an integer,

\[
\boxed{
|E|
\leq
\left\lfloor\frac{n^2}{4}\right\rfloor.
}
\]

%%%%%%%%%%%%%%%%%%%%%%%%%%%%%%%%%%%%%%%%%%%%%%%%%%%%%%%%%%%%
\subsection{Worked Example: Mantel's Theorem}
\label{subsec:worked-mantel}
%%%%%%%%%%%%%%%%%%%%%%%%%%%%%%%%%%%%%%%%%%%%%%%%%%%%%%%%%%%%

\begin{workedexamplebox}{Maximum Edges in a Triangle-Free Graph}

Suppose a graph has

\[
n=10
\]

vertices and contains no triangle.

Mantel's theorem gives

\[
|E|
\leq
\left\lfloor
\frac{10^2}{4}
\right\rfloor
=
25.
\]

The complete bipartite graph

\[
K_{5,5}
\]

has exactly \(25\) edges and is triangle-free.

\begin{answerbox}
\[
\boxed{
\operatorname{ex}(10,K_3)=25.
}
\]
\end{answerbox}

\end{workedexamplebox}

%%%%%%%%%%%%%%%%%%%%%%%%%%%%%%%%%%%%%%%%%%%%%%%%%%%%%%%%%%%%
\subsection{Tur\'an Graphs}
\label{subsec:turan-graphs}
%%%%%%%%%%%%%%%%%%%%%%%%%%%%%%%%%%%%%%%%%%%%%%%%%%%%%%%%%%%%

Mantel's theorem generalizes to larger complete graphs.

\begin{definitionbox}{Tur\'an Graph}
The Tur\'an graph \(T_r(n)\) is the complete \(r\)-partite graph on
\(n\) vertices whose part sizes differ by at most one.
\end{definitionbox}

Thus the vertices are divided as evenly as possible into \(r\) parts.

Vertices in different parts are adjacent.

Vertices in the same part are not adjacent.

%%%%%%%%%%%%%%%%%%%%%%%%%%%%%%%%%%%%%%%%%%%%%%%%%%%%%%%%%%%%
\subsection{Tur\'an's Theorem}
\label{subsec:turans-theorem}
%%%%%%%%%%%%%%%%%%%%%%%%%%%%%%%%%%%%%%%%%%%%%%%%%%%%%%%%%%%%

\begin{theorem}[Tur\'an's Theorem]
Among all \(n\)-vertex graphs containing no copy of \(K_{r+1}\), the
Tur\'an graph \(T_r(n)\) has the maximum possible number of edges.

Equivalently,

\[
\boxed{
\operatorname{ex}(n,K_{r+1})
=
|E(T_r(n))|.
}
\]
\end{theorem}

Mantel's theorem is the special case

\[
r=2.
\]

%%%%%%%%%%%%%%%%%%%%%%%%%%%%%%%%%%%%%%%%%%%%%%%%%%%%%%%%%%%%
\subsection{Approximate Size of a Tur\'an Graph}
\label{subsec:turan-graph-size}
%%%%%%%%%%%%%%%%%%%%%%%%%%%%%%%%%%%%%%%%%%%%%%%%%%%%%%%%%%%%

When the \(r\) parts are nearly equal, approximately

\[
\frac1r
\]

of all vertex pairs lie within the same part and are therefore missing
as edges.

Thus,

\[
|E(T_r(n))|
\approx
\left(
1-\frac1r
\right)
\frac{n^2}{2}.
\]

More precisely,

\[
\boxed{
\operatorname{ex}(n,K_{r+1})
=
\left(
1-\frac1r
\right)
\frac{n^2}{2}
+
O(1).
}
\]

%%%%%%%%%%%%%%%%%%%%%%%%%%%%%%%%%%%%%%%%%%%%%%%%%%%%%%%%%%%%
\subsection{Worked Example: Forbidding \(K_4\)}
\label{subsec:worked-turan-k4}
%%%%%%%%%%%%%%%%%%%%%%%%%%%%%%%%%%%%%%%%%%%%%%%%%%%%%%%%%%%%

\begin{workedexamplebox}{Maximum Edges Without a \(K_4\)}

Suppose

\[
n=9.
\]

To avoid \(K_4\), Tur\'an's theorem uses

\[
T_3(9).
\]

The vertices are divided into three equal parts of size \(3\).

Every pair of parts contributes

\[
3\cdot3=9
\]

edges.

There are

\[
\binom32=3
\]

pairs of parts.

Therefore,

\[
|E(T_3(9))|
=
3\cdot9
=
27.
\]

\begin{answerbox}
\[
\boxed{
\operatorname{ex}(9,K_4)=27.
}
\]
\end{answerbox}

\end{workedexamplebox}

%%%%%%%%%%%%%%%%%%%%%%%%%%%%%%%%%%%%%%%%%%%%%%%%%%%%%%%%%%%%
\subsection{Chromatic Number and Extremal Density}
\label{subsec:chromatic-number-extremal-density}
%%%%%%%%%%%%%%%%%%%%%%%%%%%%%%%%%%%%%%%%%%%%%%%%%%%%%%%%%%%%

For a graph \(H\), let

\[
\chi(H)
\]

denote its chromatic number.

The chromatic number turns out to determine the leading asymptotic term
of

\[
\operatorname{ex}(n,H)
\]

for every fixed non-bipartite graph \(H\).

This is the content of the Erd\H{o}s--Stone theorem.

%%%%%%%%%%%%%%%%%%%%%%%%%%%%%%%%%%%%%%%%%%%%%%%%%%%%%%%%%%%%
\subsection{Erd\H{o}s--Stone Theorem}
\label{subsec:erdos-stone}
%%%%%%%%%%%%%%%%%%%%%%%%%%%%%%%%%%%%%%%%%%%%%%%%%%%%%%%%%%%%

\begin{theorem}[Erd\H{o}s--Stone Theorem]
For a fixed graph \(H\) with

\[
\chi(H)=r+1,
\]

we have

\[
\boxed{
\operatorname{ex}(n,H)
=
\left(
1-\frac1r+o(1)
\right)
\frac{n^2}{2}.
}
\]

Here \(o(1)\) denotes a quantity tending to \(0\) as \(n\to\infty\).
\end{theorem}

Thus, asymptotically, the chromatic number of \(H\) determines the
maximum possible edge density of an \(H\)-free graph.

%%%%%%%%%%%%%%%%%%%%%%%%%%%%%%%%%%%%%%%%%%%%%%%%%%%%%%%%%%%%
\subsection{The Bipartite Exception}
\label{subsec:bipartite-extremal-exception}
%%%%%%%%%%%%%%%%%%%%%%%%%%%%%%%%%%%%%%%%%%%%%%%%%%%%%%%%%%%%

If \(H\) is bipartite, then

\[
\chi(H)=2.
\]

The Erd\H{o}s--Stone leading density becomes zero.

Thus,

\[
\operatorname{ex}(n,H)=o(n^2).
\]

Determining the correct smaller-order growth can be much harder.

For example, forbidding even cycles or complete bipartite graphs leads
to important problems in extremal graph theory.

%%%%%%%%%%%%%%%%%%%%%%%%%%%%%%%%%%%%%%%%%%%%%%%%%%%%%%%%%%%%
\subsection{Set Systems and Extremal Questions}
\label{subsec:set-systems-extremal}
%%%%%%%%%%%%%%%%%%%%%%%%%%%%%%%%%%%%%%%%%%%%%%%%%%%%%%%%%%%%

Extremal combinatorics is not limited to graphs.

Let

\[
[n]
=
\{1,2,\ldots,n\}.
\]

A family

\[
\mathcal{F}
\subseteq
\mathcal{P}([n])
\]

may be restricted by intersection or containment conditions.

One asks how large \(\mathcal{F}\) can be.

These questions form extremal set theory.

%%%%%%%%%%%%%%%%%%%%%%%%%%%%%%%%%%%%%%%%%%%%%%%%%%%%%%%%%%%%
\subsection{Intersecting Families}
\label{subsec:intersecting-families}
%%%%%%%%%%%%%%%%%%%%%%%%%%%%%%%%%%%%%%%%%%%%%%%%%%%%%%%%%%%%

\begin{definitionbox}{Intersecting Family}
A family of sets \(\mathcal{F}\) is intersecting if

\[
\boxed{
A\cap B\neq\varnothing
}
\]

for every

\[
A,B\in\mathcal{F}.
\]
\end{definitionbox}

A natural question is:

\[
\boxed{
\text{How large can an intersecting family of }k\text{-subsets of }
[n]\text{ be?}
}
\]

The answer is given by the Erd\H{o}s--Ko--Rado theorem.

%%%%%%%%%%%%%%%%%%%%%%%%%%%%%%%%%%%%%%%%%%%%%%%%%%%%%%%%%%%%
\subsection{Erd\H{o}s--Ko--Rado Theorem}
\label{subsec:erdos-ko-rado}
%%%%%%%%%%%%%%%%%%%%%%%%%%%%%%%%%%%%%%%%%%%%%%%%%%%%%%%%%%%%

\begin{theorem}[Erd\H{o}s--Ko--Rado Theorem]
Suppose

\[
n\geq2k.
\]

If

\[
\mathcal{F}
\subseteq
\binom{[n]}{k}
\]

is intersecting, then

\[
\boxed{
|\mathcal{F}|
\leq
\binom{n-1}{k-1}.
}
\]

For \(n>2k\), equality is achieved by taking all \(k\)-subsets containing
one fixed element.
\end{theorem}

The standard extremal family is therefore

\[
\boxed{
\mathcal{F}
=
\left\{
A\in\binom{[n]}{k}:1\in A
\right\}.
}
\]

%%%%%%%%%%%%%%%%%%%%%%%%%%%%%%%%%%%%%%%%%%%%%%%%%%%%%%%%%%%%
\subsection{Worked Example: Intersecting Families}
\label{subsec:worked-ekr}
%%%%%%%%%%%%%%%%%%%%%%%%%%%%%%%%%%%%%%%%%%%%%%%%%%%%%%%%%%%%

\begin{workedexamplebox}{Intersecting \(3\)-Subsets of \([8]\)}

Let

\[
\mathcal{F}
\subseteq
\binom{[8]}{3}
\]

be intersecting.

Since

\[
8\geq2(3),
\]

the Erd\H{o}s--Ko--Rado theorem applies.

Therefore,

\[
|\mathcal{F}|
\leq
\binom{7}{2}.
\]

\begin{answerbox}
\[
\boxed{
|\mathcal{F}|\leq21.
}
\]
\end{answerbox}

\end{workedexamplebox}

%%%%%%%%%%%%%%%%%%%%%%%%%%%%%%%%%%%%%%%%%%%%%%%%%%%%%%%%%%%%
\subsection{Stars in Set Systems}
\label{subsec:stars-set-systems}
%%%%%%%%%%%%%%%%%%%%%%%%%%%%%%%%%%%%%%%%%%%%%%%%%%%%%%%%%%%%

A family of the form

\[
\mathcal{F}_i
=
\left\{
A\in\binom{[n]}{k}:i\in A
\right\}
\]

is called a \emph{star}.

Its size is

\[
\boxed{
|\mathcal{F}_i|
=
\binom{n-1}{k-1}.
}
\]

The Erd\H{o}s--Ko--Rado theorem states that, under the appropriate
conditions, no intersecting \(k\)-uniform family can be larger.

%%%%%%%%%%%%%%%%%%%%%%%%%%%%%%%%%%%%%%%%%%%%%%%%%%%%%%%%%%%%
\subsection{The Boolean Lattice}
\label{subsec:boolean-lattice-extremal}
%%%%%%%%%%%%%%%%%%%%%%%%%%%%%%%%%%%%%%%%%%%%%%%%%%%%%%%%%%%%

The power set

\[
\mathcal{P}([n])
\]

ordered by inclusion is called the Boolean lattice.

Its levels are

\[
\boxed{
\binom{[n]}{0},
\binom{[n]}{1},
\ldots,
\binom{[n]}{n}.
}
\]

The \(k\)-th level contains

\[
\binom{n}{k}
\]

sets.

The largest level occurs near

\[
k=\frac n2.
\]

%%%%%%%%%%%%%%%%%%%%%%%%%%%%%%%%%%%%%%%%%%%%%%%%%%%%%%%%%%%%
\subsection{Antichains}
\label{subsec:antichains-extremal}
%%%%%%%%%%%%%%%%%%%%%%%%%%%%%%%%%%%%%%%%%%%%%%%%%%%%%%%%%%%%

An antichain is a family of sets in which no member contains another.

Thus,

\[
A,B\in\mathcal{F},
\qquad
A\neq B
\]

implies

\[
A\nsubseteq B
\]

and

\[
B\nsubseteq A.
\]

The natural extremal question is:

\[
\boxed{
\text{What is the largest antichain in }\mathcal{P}([n])?
}
\]

%%%%%%%%%%%%%%%%%%%%%%%%%%%%%%%%%%%%%%%%%%%%%%%%%%%%%%%%%%%%
\subsection{Sperner's Theorem}
\label{subsec:sperners-theorem}
%%%%%%%%%%%%%%%%%%%%%%%%%%%%%%%%%%%%%%%%%%%%%%%%%%%%%%%%%%%%

\begin{theorem}[Sperner's Theorem]
If

\[
\mathcal{F}
\subseteq
\mathcal{P}([n])
\]

is an antichain, then

\[
\boxed{
|\mathcal{F}|
\leq
\binom{n}{\lfloor n/2\rfloor}.
}
\]

Equality is attained by a middle level of the Boolean lattice.
\end{theorem}

Thus the largest possible antichain consists of all subsets of one
middle size.

%%%%%%%%%%%%%%%%%%%%%%%%%%%%%%%%%%%%%%%%%%%%%%%%%%%%%%%%%%%%
\subsection{Worked Example: Sperner's Theorem}
\label{subsec:worked-sperner}
%%%%%%%%%%%%%%%%%%%%%%%%%%%%%%%%%%%%%%%%%%%%%%%%%%%%%%%%%%%%

\begin{workedexamplebox}{Largest Antichain in \(\mathcal{P}([6])\)}

For

\[
n=6,
\]

the middle level has size

\[
\binom63=20.
\]

Therefore any antichain satisfies

\[
|\mathcal{F}|\leq20.
\]

The family of all \(3\)-element subsets is itself an antichain.

\begin{answerbox}
\[
\boxed{
\max|\mathcal{F}|=20.
}
\]
\end{answerbox}

\end{workedexamplebox}

%%%%%%%%%%%%%%%%%%%%%%%%%%%%%%%%%%%%%%%%%%%%%%%%%%%%%%%%%%%%
\subsection{Chains in the Boolean Lattice}
\label{subsec:chains-boolean-lattice}
%%%%%%%%%%%%%%%%%%%%%%%%%%%%%%%%%%%%%%%%%%%%%%%%%%%%%%%%%%%%

A maximal chain in

\[
\mathcal{P}([n])
\]

has the form

\[
\varnothing
\subset
A_1
\subset
A_2
\subset
\cdots
\subset
A_n=[n],
\]

where

\[
|A_k|=k.
\]

Each maximal chain corresponds to an ordering of the \(n\) elements.

Therefore, the Boolean lattice contains

\[
\boxed{
n!
}
\]

maximal chains.

%%%%%%%%%%%%%%%%%%%%%%%%%%%%%%%%%%%%%%%%%%%%%%%%%%%%%%%%%%%%
\subsection{Lubell Function}
\label{subsec:lubell-function}
%%%%%%%%%%%%%%%%%%%%%%%%%%%%%%%%%%%%%%%%%%%%%%%%%%%%%%%%%%%%

For a family

\[
\mathcal{F}
\subseteq
\mathcal{P}([n]),
\]

define its Lubell function by

\[
\boxed{
L(\mathcal{F})
=
\sum_{A\in\mathcal{F}}
\frac{1}{\binom{n}{|A|}}.
}
\]

This quantity measures the expected number of members of
\(\mathcal{F}\) lying on a uniformly random maximal chain.

%%%%%%%%%%%%%%%%%%%%%%%%%%%%%%%%%%%%%%%%%%%%%%%%%%%%%%%%%%%%
\subsection{LYM Inequality}
\label{subsec:lym-inequality}
%%%%%%%%%%%%%%%%%%%%%%%%%%%%%%%%%%%%%%%%%%%%%%%%%%%%%%%%%%%%

\begin{theorem}[Lubell--Yamamoto--Meshalkin Inequality]
If \(\mathcal{F}\) is an antichain in \(\mathcal{P}([n])\), then

\[
\boxed{
\sum_{A\in\mathcal{F}}
\frac{1}{\binom{n}{|A|}}
\leq1.
}
\]
\end{theorem}

This is commonly called the LYM inequality.

Since

\[
\binom{n}{|A|}
\leq
\binom{n}{\lfloor n/2\rfloor},
\]

we obtain

\[
\frac{|\mathcal{F}|}{
\binom{n}{\lfloor n/2\rfloor}
}
\leq1,
\]

which yields Sperner's theorem.

%%%%%%%%%%%%%%%%%%%%%%%%%%%%%%%%%%%%%%%%%%%%%%%%%%%%%%%%%%%%
\subsection{Double Counting Proof of the LYM Inequality}
\label{subsec:lym-double-counting}
%%%%%%%%%%%%%%%%%%%%%%%%%%%%%%%%%%%%%%%%%%%%%%%%%%%%%%%%%%%%

Count pairs

\[
(A,\mathcal{C}),
\]

where

\[
A\in\mathcal{F}
\]

and \(\mathcal{C}\) is a maximal chain containing \(A\).

If

\[
|A|=k,
\]

then the number of maximal chains through \(A\) is

\[
k!(n-k)!.
\]

Therefore, the total number of such pairs is

\[
\sum_{A\in\mathcal{F}}
|A|!(n-|A|)!.
\]

Since \(\mathcal{F}\) is an antichain, each maximal chain contains at
most one member of \(\mathcal{F}\).

There are \(n!\) maximal chains.

Hence,

\[
\sum_{A\in\mathcal{F}}
|A|!(n-|A|)!
\leq
n!.
\]

Divide by \(n!\):

\[
\boxed{
\sum_{A\in\mathcal{F}}
\frac{1}{\binom{n}{|A|}}
\leq1.
}
\]

%%%%%%%%%%%%%%%%%%%%%%%%%%%%%%%%%%%%%%%%%%%%%%%%%%%%%%%%%%%%
\subsection{Dilworth's Theorem}
\label{subsec:dilworth-extremal}
%%%%%%%%%%%%%%%%%%%%%%%%%%%%%%%%%%%%%%%%%%%%%%%%%%%%%%%%%%%%

Extremal questions about chains and antichains extend to arbitrary
finite partially ordered sets.

\begin{theorem}[Dilworth's Theorem]
In every finite poset,

\[
\boxed{
\text{maximum size of an antichain}
=
\text{minimum number of chains needed to partition the poset}.
}
\]
\end{theorem}

This is another example of a min-max theorem in discrete mathematics.

%%%%%%%%%%%%%%%%%%%%%%%%%%%%%%%%%%%%%%%%%%%%%%%%%%%%%%%%%%%%
\subsection{Mirsky's Theorem}
\label{subsec:mirsky-extremal}
%%%%%%%%%%%%%%%%%%%%%%%%%%%%%%%%%%%%%%%%%%%%%%%%%%%%%%%%%%%%

The dual statement is Mirsky's theorem.

\begin{theorem}[Mirsky's Theorem]
In every finite poset,

\[
\boxed{
\text{maximum size of a chain}
=
\text{minimum number of antichains needed to partition the poset}.
}
\]
\end{theorem}

Dilworth's and Mirsky's theorems reveal deep structural duality between
chains and antichains.

%%%%%%%%%%%%%%%%%%%%%%%%%%%%%%%%%%%%%%%%%%%%%%%%%%%%%%%%%%%%
\subsection{Compression and Shifting}
\label{subsec:compression-shifting}
%%%%%%%%%%%%%%%%%%%%%%%%%%%%%%%%%%%%%%%%%%%%%%%%%%%%%%%%%%%%

Compression methods transform a set family into a more structured
family without changing its size and while preserving relevant
properties.

A basic shift attempts to replace a larger element by a smaller one.

For \(i<j\), define a shift that replaces \(j\) by \(i\) when doing so
creates a new admissible set.

Repeated shifting often moves the family toward a canonical extremal
shape.

%%%%%%%%%%%%%%%%%%%%%%%%%%%%%%%%%%%%%%%%%%%%%%%%%%%%%%%%%%%%
\subsection{Why Compression Is Useful}
\label{subsec:why-compression-useful}
%%%%%%%%%%%%%%%%%%%%%%%%%%%%%%%%%%%%%%%%%%%%%%%%%%%%%%%%%%%%

An arbitrary extremal family may be difficult to analyze.

Compression can simplify the family while preserving:

\begin{itemize}
    \item cardinality;
    \item uniformity;
    \item intersection properties;
    \item certain forbidden configurations.
\end{itemize}

The resulting compressed family is often easier to classify.

Compression methods are important in proofs of results such as the
Erd\H{o}s--Ko--Rado theorem.

%%%%%%%%%%%%%%%%%%%%%%%%%%%%%%%%%%%%%%%%%%%%%%%%%%%%%%%%%%%%
\subsection{Averaging Arguments}
\label{subsec:averaging-extremal}
%%%%%%%%%%%%%%%%%%%%%%%%%%%%%%%%%%%%%%%%%%%%%%%%%%%%%%%%%%%%

Averaging is a common extremal technique.

Suppose quantities

\[
x_1,\ldots,x_n
\]

have average

\[
\overline{x}
=
\frac1n
\sum_{i=1}^{n}x_i.
\]

Then at least one value satisfies

\[
x_i\geq\overline{x}
\]

and at least one satisfies

\[
x_j\leq\overline{x}.
\]

Simple as this principle is, it can identify vertices, subsets, or local
configurations with unusually large or small structure.

%%%%%%%%%%%%%%%%%%%%%%%%%%%%%%%%%%%%%%%%%%%%%%%%%%%%%%%%%%%%
\subsection{Average Degree}
\label{subsec:average-degree-extremal}
%%%%%%%%%%%%%%%%%%%%%%%%%%%%%%%%%%%%%%%%%%%%%%%%%%%%%%%%%%%%

For a graph

\[
G=(V,E),
\]

the average degree is

\[
\boxed{
\overline{d}
=
\frac{2|E|}{|V|}.
}
\]

Therefore some vertex has degree at least

\[
\frac{2|E|}{|V|},
\]

and some vertex has degree at most this quantity.

Average-degree arguments frequently appear in extremal graph proofs.

%%%%%%%%%%%%%%%%%%%%%%%%%%%%%%%%%%%%%%%%%%%%%%%%%%%%%%%%%%%%
\subsection{Double Counting in Extremal Combinatorics}
\label{subsec:double-counting-extremal}
%%%%%%%%%%%%%%%%%%%%%%%%%%%%%%%%%%%%%%%%%%%%%%%%%%%%%%%%%%%%

Double counting often compares local and global structure.

One chooses an incidence set such as

\[
\boxed{
\{
(x,Y):
x\text{ is incident with }Y
\}
}
\]

and counts it in two ways.

Examples include:

\begin{itemize}
    \item vertex-edge incidences;
    \item set-chain incidences;
    \item vertex-clique incidences;
    \item point-block incidences;
    \item edge-triangle incidences.
\end{itemize}

Many extremal inequalities arise from such comparisons.

%%%%%%%%%%%%%%%%%%%%%%%%%%%%%%%%%%%%%%%%%%%%%%%%%%%%%%%%%%%%
\subsection{The Probabilistic Method}
\label{subsec:probabilistic-method-extremal}
%%%%%%%%%%%%%%%%%%%%%%%%%%%%%%%%%%%%%%%%%%%%%%%%%%%%%%%%%%%%

A striking technique in extremal combinatorics proves existence using
randomness.

The basic principle is:

\[
\boxed{
\mathbb{P}(\text{desired property})>0
\Longrightarrow
\text{at least one desired object exists}.
}
\]

One does not need to identify the object explicitly.

Instead, one proves that a random object has the desired property with
positive probability.

%%%%%%%%%%%%%%%%%%%%%%%%%%%%%%%%%%%%%%%%%%%%%%%%%%%%%%%%%%%%
\subsection{Expectation Method}
\label{subsec:expectation-method}
%%%%%%%%%%%%%%%%%%%%%%%%%%%%%%%%%%%%%%%%%%%%%%%%%%%%%%%%%%%%

Suppose a random variable \(X\) measures some undesirable quantity.

If

\[
\mathbb{E}[X]<1
\]

and \(X\) is a nonnegative integer-valued random variable, then some
outcome must satisfy

\[
X=0.
\]

Otherwise \(X\geq1\) always, forcing

\[
\mathbb{E}[X]\geq1.
\]

Thus expectation alone can prove existence.

%%%%%%%%%%%%%%%%%%%%%%%%%%%%%%%%%%%%%%%%%%%%%%%%%%%%%%%%%%%%
\subsection{Worked Example: Independent Sets by Random Selection}
\label{subsec:worked-probabilistic-independent-set}
%%%%%%%%%%%%%%%%%%%%%%%%%%%%%%%%%%%%%%%%%%%%%%%%%%%%%%%%%%%%

\begin{workedexamplebox}{A Simple Random-Subset Argument}

Let \(G\) have \(n\) vertices and \(m\) edges.

Choose each vertex independently with probability \(p\).

Let \(X\) be the number of selected vertices and \(Y\) the number of
edges having both endpoints selected.

Then

\[
\mathbb{E}[X]=pn
\]

and

\[
\mathbb{E}[Y]=p^2m.
\]

Delete one endpoint from every selected edge.

The remaining set is independent and has size at least

\[
X-Y.
\]

Therefore some independent set has size at least

\[
pn-p^2m.
\]

\begin{answerbox}
\[
\boxed{
\alpha(G)
\geq
pn-p^2m
}
\]

for every \(0\leq p\leq1\), where \(\alpha(G)\) is the independence
number.
\end{answerbox}

\end{workedexamplebox}

%%%%%%%%%%%%%%%%%%%%%%%%%%%%%%%%%%%%%%%%%%%%%%%%%%%%%%%%%%%%
\subsection{Optimizing the Probabilistic Bound}
\label{subsec:optimizing-probabilistic-bound}
%%%%%%%%%%%%%%%%%%%%%%%%%%%%%%%%%%%%%%%%%%%%%%%%%%%%%%%%%%%%

The expression

\[
pn-p^2m
\]

is a quadratic function of \(p\).

When

\[
m>0,
\]

its maximum occurs at

\[
p=\frac{n}{2m},
\]

provided this lies in the interval \([0,1]\).

Thus probability produces a lower bound, and elementary optimization
can sharpen it.

This interaction between randomness and optimization is common in
modern combinatorics.

%%%%%%%%%%%%%%%%%%%%%%%%%%%%%%%%%%%%%%%%%%%%%%%%%%%%%%%%%%%%
\subsection{Alteration Method}
\label{subsec:alteration-method}
%%%%%%%%%%%%%%%%%%%%%%%%%%%%%%%%%%%%%%%%%%%%%%%%%%%%%%%%%%%%

The previous argument illustrates the alteration method:

\begin{enumerate}
    \item construct a random object;
    \item allow some violations;
    \item modify the object to remove the violations;
    \item show that enough structure remains.
\end{enumerate}

Random construction followed by deterministic repair is a powerful
extremal technique.

%%%%%%%%%%%%%%%%%%%%%%%%%%%%%%%%%%%%%%%%%%%%%%%%%%%%%%%%%%%%
\subsection{Random Graphs and Extremal Existence}
\label{subsec:random-graphs-extremal-preview}
%%%%%%%%%%%%%%%%%%%%%%%%%%%%%%%%%%%%%%%%%%%%%%%%%%%%%%%%%%%%

Random graphs can prove the existence of graphs satisfying seemingly
conflicting properties.

For example, one may seek graphs with:

\begin{itemize}
    \item high chromatic number;
    \item large girth;
    \item few forbidden subgraphs;
    \item small independence number.
\end{itemize}

Random construction can show that such graphs exist even when explicit
constructions are difficult.

This viewpoint will be developed further in probabilistic
combinatorics.

%%%%%%%%%%%%%%%%%%%%%%%%%%%%%%%%%%%%%%%%%%%%%%%%%%%%%%%%%%%%
\subsection{Ramsey-Theoretic Extremal Behavior}
\label{subsec:ramsey-extremal-preview}
%%%%%%%%%%%%%%%%%%%%%%%%%%%%%%%%%%%%%%%%%%%%%%%%%%%%%%%%%%%%

Ramsey theory asks when large structures must contain a prescribed
ordered or monochromatic substructure.

Extremal and Ramsey questions are closely related.

For example:

\[
\boxed{
\text{extremal theory}
:
\text{avoid a configuration as long as possible}
}
\]

while

\[
\boxed{
\text{Ramsey theory}
:
\text{determine when avoidance becomes impossible}.
}
\]

Ramsey theory will appear later in this chapter.

%%%%%%%%%%%%%%%%%%%%%%%%%%%%%%%%%%%%%%%%%%%%%%%%%%%%%%%%%%%%
\subsection{Hypergraph Extremal Problems}
\label{subsec:hypergraph-extremal-problems}
%%%%%%%%%%%%%%%%%%%%%%%%%%%%%%%%%%%%%%%%%%%%%%%%%%%%%%%%%%%%

For a \(k\)-uniform hypergraph, one may define extremal numbers just as
in graph theory.

If \(F\) is a fixed \(k\)-uniform hypergraph, define

\[
\boxed{
\operatorname{ex}_k(n,F)
}
\]

as the maximum number of hyperedges in an \(n\)-vertex \(k\)-uniform
hypergraph containing no copy of \(F\).

Hypergraph extremal problems are generally more difficult than graph
extremal problems.

Even natural analogues of classical graph questions can remain open.

%%%%%%%%%%%%%%%%%%%%%%%%%%%%%%%%%%%%%%%%%%%%%%%%%%%%%%%%%%%%
\subsection{Extremal Set Intersection Problems}
\label{subsec:extremal-intersection-problems}
%%%%%%%%%%%%%%%%%%%%%%%%%%%%%%%%%%%%%%%%%%%%%%%%%%%%%%%%%%%%

One may impose stronger intersection restrictions.

For example:

\[
|A\cap B|\geq t
\]

for every pair

\[
A,B\in\mathcal{F}.
\]

Such a family is called \(t\)-intersecting.

The question becomes:

\[
\boxed{
\text{How large can a }t\text{-intersecting family be?}
}
\]

These problems generalize the Erd\H{o}s--Ko--Rado setting.

%%%%%%%%%%%%%%%%%%%%%%%%%%%%%%%%%%%%%%%%%%%%%%%%%%%%%%%%%%%%
\subsection{Forbidden Intersection Sizes}
\label{subsec:forbidden-intersection-sizes}
%%%%%%%%%%%%%%%%%%%%%%%%%%%%%%%%%%%%%%%%%%%%%%%%%%%%%%%%%%%%

Another family of extremal problems forbids particular intersection
sizes.

For example, require

\[
|A\cap B|\neq r
\]

for distinct members \(A,B\).

Such restrictions connect extremal set theory with:

\begin{itemize}
    \item linear algebra;
    \item coding theory;
    \item finite geometry;
    \item polynomial methods.
\end{itemize}

%%%%%%%%%%%%%%%%%%%%%%%%%%%%%%%%%%%%%%%%%%%%%%%%%%%%%%%%%%%%
\subsection{Linear Algebra Method}
\label{subsec:linear-algebra-method-extremal}
%%%%%%%%%%%%%%%%%%%%%%%%%%%%%%%%%%%%%%%%%%%%%%%%%%%%%%%%%%%%

Linear algebra can provide extremal bounds by associating vectors or
polynomials with combinatorial objects.

If too many objects existed, the associated vectors would violate a
dimension bound or become linearly dependent in an impossible way.

The general principle is

\[
\boxed{
\text{combinatorial family}
\longrightarrow
\text{vectors or polynomials}
\longrightarrow
\text{dimension bound}.
}
\]

This method is especially powerful in intersection problems.

%%%%%%%%%%%%%%%%%%%%%%%%%%%%%%%%%%%%%%%%%%%%%%%%%%%%%%%%%%%%
\subsection{Polynomial Method}
\label{subsec:polynomial-method-extremal}
%%%%%%%%%%%%%%%%%%%%%%%%%%%%%%%%%%%%%%%%%%%%%%%%%%%%%%%%%%%%

The polynomial method encodes combinatorial restrictions into
polynomials.

Typical arguments involve:

\begin{itemize}
    \item interpolation;
    \item polynomial vanishing;
    \item degree bounds;
    \item linear independence of polynomial families;
    \item finite-field identities.
\end{itemize}

The method has become a major tool in modern combinatorics and finite
geometry.

%%%%%%%%%%%%%%%%%%%%%%%%%%%%%%%%%%%%%%%%%%%%%%%%%%%%%%%%%%%%
\subsection{Stability}
\label{subsec:stability-extremal}
%%%%%%%%%%%%%%%%%%%%%%%%%%%%%%%%%%%%%%%%%%%%%%%%%%%%%%%%%%%%

An exact extremal theorem identifies the largest possible structure.

A stability theorem asks:

\[
\boxed{
\text{If a structure is nearly extremal, must it resemble an extremal
example?}
}
\]

For many graph problems, the answer is yes.

For example, a triangle-free graph with nearly

\[
\frac{n^2}{4}
\]

edges must be structurally close to a complete bipartite graph.

%%%%%%%%%%%%%%%%%%%%%%%%%%%%%%%%%%%%%%%%%%%%%%%%%%%%%%%%%%%%
\subsection{Why Stability Matters}
\label{subsec:why-stability-matters}
%%%%%%%%%%%%%%%%%%%%%%%%%%%%%%%%%%%%%%%%%%%%%%%%%%%%%%%%%%%%

An extremal theorem gives the maximum size.

A stability theorem gives structural information about all
near-maximizers.

Thus stability strengthens

\[
\boxed{
\text{optimal value}
}
\]

into

\[
\boxed{
\text{optimal value}
+
\text{shape of nearly optimal objects}.
}
\]

This is important in asymptotic combinatorics and probabilistic
structure theory.

%%%%%%%%%%%%%%%%%%%%%%%%%%%%%%%%%%%%%%%%%%%%%%%%%%%%%%%%%%%%
\subsection{Supersaturation}
\label{subsec:supersaturation}
%%%%%%%%%%%%%%%%%%%%%%%%%%%%%%%%%%%%%%%%%%%%%%%%%%%%%%%%%%%%

Suppose an \(n\)-vertex graph has more edges than the extremal threshold
for avoiding \(H\).

Then at least one copy of \(H\) must appear.

Supersaturation asks for more:

\[
\boxed{
\text{How many copies of }H\text{ must appear once the threshold is
exceeded?}
}
\]

Often, adding substantially more than

\[
\operatorname{ex}(n,H)
\]

edges forces many copies of \(H\).

%%%%%%%%%%%%%%%%%%%%%%%%%%%%%%%%%%%%%%%%%%%%%%%%%%%%%%%%%%%%
\subsection{Removal Lemmas}
\label{subsec:removal-lemmas-preview}
%%%%%%%%%%%%%%%%%%%%%%%%%%%%%%%%%%%%%%%%%%%%%%%%%%%%%%%%%%%%

A removal lemma has the general flavor:

\[
\boxed{
\text{few forbidden configurations}
\Longrightarrow
\text{small modification removes all of them}.
}
\]

For example, the triangle removal lemma says roughly that if a large
graph contains relatively few triangles, then deleting relatively few
edges can make the graph triangle-free.

Removal lemmas connect extremal combinatorics, property testing, and
additive combinatorics.

%%%%%%%%%%%%%%%%%%%%%%%%%%%%%%%%%%%%%%%%%%%%%%%%%%%%%%%%%%%%
\subsection{Exact Versus Asymptotic Extremal Results}
\label{subsec:exact-asymptotic-extremal}
%%%%%%%%%%%%%%%%%%%%%%%%%%%%%%%%%%%%%%%%%%%%%%%%%%%%%%%%%%%%

An exact result determines a precise value such as

\[
\operatorname{ex}(n,K_3)
=
\left\lfloor\frac{n^2}{4}\right\rfloor.
\]

An asymptotic result may only determine leading growth:

\[
\operatorname{ex}(n,H)
=
\left(
1-\frac{1}{\chi(H)-1}+o(1)
\right)
\frac{n^2}{2}.
\]

Both are important.

Exact results reveal detailed finite structure.

Asymptotic results reveal large-scale behavior.

%%%%%%%%%%%%%%%%%%%%%%%%%%%%%%%%%%%%%%%%%%%%%%%%%%%%%%%%%%%%
\subsection{Extremal Configurations}
\label{subsec:extremal-configurations}
%%%%%%%%%%%%%%%%%%%%%%%%%%%%%%%%%%%%%%%%%%%%%%%%%%%%%%%%%%%%

An extremal configuration is a structure attaining the maximum or
minimum allowed value.

Examples include:

\[
K_{\lfloor n/2\rfloor,\lceil n/2\rceil}
\]

for Mantel's theorem and

\[
T_r(n)
\]

for Tur\'an's theorem.

In set systems, a full star is extremal in the classical
Erd\H{o}s--Ko--Rado setting.

A middle level of the Boolean lattice is extremal in Sperner's theorem.

%%%%%%%%%%%%%%%%%%%%%%%%%%%%%%%%%%%%%%%%%%%%%%%%%%%%%%%%%%%%
\subsection{Uniqueness of Extremal Configurations}
\label{subsec:uniqueness-extremal}
%%%%%%%%%%%%%%%%%%%%%%%%%%%%%%%%%%%%%%%%%%%%%%%%%%%%%%%%%%%%

An extremal value may be attained by:

\begin{itemize}
    \item one structure up to isomorphism;
    \item several structurally different examples;
    \item a large family of extremal configurations.
\end{itemize}

Thus one may ask both

\[
\boxed{
\text{What is the extremal value?}
}
\]

and

\[
\boxed{
\text{Which objects attain it?}
}
\]

The second question is often substantially harder.

%%%%%%%%%%%%%%%%%%%%%%%%%%%%%%%%%%%%%%%%%%%%%%%%%%%%%%%%%%%%
\subsection{Weighted Extremal Problems}
\label{subsec:weighted-extremal-problems}
%%%%%%%%%%%%%%%%%%%%%%%%%%%%%%%%%%%%%%%%%%%%%%%%%%%%%%%%%%%%

Instead of maximizing cardinality, one may assign weights.

If each element \(x\) has weight \(w(x)\), one seeks

\[
\boxed{
\max_{A\in\mathcal{F}}
\sum_{x\in A}w(x).
}
\]

This creates a bridge between extremal combinatorics and combinatorial
optimization.

Weighted versions of matching, independent-set, and set-family
problems are especially important in applications.

%%%%%%%%%%%%%%%%%%%%%%%%%%%%%%%%%%%%%%%%%%%%%%%%%%%%%%%%%%%%
\subsection{Applications}
\label{subsec:extremal-combinatorics-applications}
%%%%%%%%%%%%%%%%%%%%%%%%%%%%%%%%%%%%%%%%%%%%%%%%%%%%%%%%%%%%

\begin{applicationbox}

\textbf{Graph Theory.}
Extremal methods determine how many edges a graph may contain while
avoiding specified subgraphs.

\medskip

\textbf{Coding Theory.}
Large collections of codewords with restricted pairwise intersections
or distances are extremal set systems.

\medskip

\textbf{Design Theory.}
Block designs impose strong intersection and incidence constraints that
naturally lead to extremal bounds.

\medskip

\textbf{Computer Science.}
Extremal bounds appear in complexity theory, property testing,
communication networks, and algorithm design.

\medskip

\textbf{Discrete Geometry.}
Questions about incidences, distances, and forbidden configurations are
often extremal.

\medskip

\textbf{Additive Combinatorics.}
Large sets avoiding arithmetic configurations form a major class of
extremal problems.

\medskip

\textbf{Probability.}
The probabilistic method proves the existence of large structures with
specified extremal properties.

\medskip

\textbf{Optimization.}
Extremal problems can often be viewed as discrete optimization problems
with highly structured feasible sets.

\end{applicationbox}

%%%%%%%%%%%%%%%%%%%%%%%%%%%%%%%%%%%%%%%%%%%%%%%%%%%%%%%%%%%%
\subsection{Common Errors}
\label{subsec:extremal-combinatorics-common-errors}
%%%%%%%%%%%%%%%%%%%%%%%%%%%%%%%%%%%%%%%%%%%%%%%%%%%%%%%%%%%%

\subsubsection{Confusing an Extremal Bound with a Construction}

An upper bound shows that no structure can be larger.

A construction shows that a structure of a certain size exists.

To prove an exact extremal result, both are usually needed.

%%%%%%%%%%%%%%%%%%%%%%%%%%%%%%%%%%%%%%%%%%%%%%%%%%%%%%%%%%%%
\subsubsection{Ignoring Sharpness}

A bound

\[
|X|\leq M
\]

is sharp only if some valid example satisfies

\[
|X|=M.
\]

%%%%%%%%%%%%%%%%%%%%%%%%%%%%%%%%%%%%%%%%%%%%%%%%%%%%%%%%%%%%
\subsubsection{Confusing \(H\)-Free with Induced-\(H\)-Free}

An \(H\)-free graph contains no subgraph isomorphic to \(H\).

An induced-\(H\)-free graph contains no induced subgraph isomorphic to
\(H\).

These are different conditions.

%%%%%%%%%%%%%%%%%%%%%%%%%%%%%%%%%%%%%%%%%%%%%%%%%%%%%%%%%%%%
\subsubsection{Applying Erd\H{o}s--Ko--Rado Without Checking \(n\geq2k\)}

The classical theorem requires

\[
n\geq2k.
\]

Ignoring this condition can lead to incorrect conclusions.

%%%%%%%%%%%%%%%%%%%%%%%%%%%%%%%%%%%%%%%%%%%%%%%%%%%%%%%%%%%%
\subsubsection{Confusing Intersecting with Pairwise Distinct}

An intersecting family requires

\[
A\cap B\neq\varnothing
\]

for every pair of members.

Distinctness alone says nothing about intersection.

%%%%%%%%%%%%%%%%%%%%%%%%%%%%%%%%%%%%%%%%%%%%%%%%%%%%%%%%%%%%
\subsubsection{Confusing Chains and Antichains}

A chain requires comparability.

An antichain requires incomparability.

These are opposite structural conditions.

%%%%%%%%%%%%%%%%%%%%%%%%%%%%%%%%%%%%%%%%%%%%%%%%%%%%%%%%%%%%
\subsubsection{Using an Asymptotic Formula as an Exact Equality}

The expression

\[
f(n)
=
g(n)+o(n^2)
\]

does not mean

\[
f(n)=g(n)
\]

for every finite \(n\).

%%%%%%%%%%%%%%%%%%%%%%%%%%%%%%%%%%%%%%%%%%%%%%%%%%%%%%%%%%%%
\subsubsection{Assuming a Random Construction Gives an Explicit Object}

The probabilistic method may prove existence without identifying a
specific example.

Existence and construction are different questions.

%%%%%%%%%%%%%%%%%%%%%%%%%%%%%%%%%%%%%%%%%%%%%%%%%%%%%%%%%%%%
\subsubsection{Ignoring Near-Extremal Structure}

Knowing the extremal value does not automatically describe structures
whose size is close to that value.

This is the role of stability theory.

%%%%%%%%%%%%%%%%%%%%%%%%%%%%%%%%%%%%%%%%%%%%%%%%%%%%%%%%%%%%
\subsection{Exercises}
\label{subsec:extremal-combinatorics-exercises}
%%%%%%%%%%%%%%%%%%%%%%%%%%%%%%%%%%%%%%%%%%%%%%%%%%%%%%%%%%%%

\begin{exercise}[1]
Define

\[
\operatorname{ex}(n,H).
\]
\end{exercise}

\begin{exercise}[1]
State Mantel's theorem.
\end{exercise}

\begin{exercise}[1]
Compute

\[
\operatorname{ex}(8,K_3).
\]
\end{exercise}

\begin{exercise}[1]
Give a triangle-free graph on \(8\) vertices having the maximum possible
number of edges.
\end{exercise}

\begin{exercise}[1]
Define the Tur\'an graph \(T_r(n)\).
\end{exercise}

\begin{exercise}[1]
Determine the part sizes of

\[
T_3(10).
\]
\end{exercise}

\begin{exercise}[1]
Define an intersecting family of sets.
\end{exercise}

\begin{exercise}[1]
Define an antichain.
\end{exercise}

\begin{exercise}[1]
Find the size of the largest antichain in

\[
\mathcal{P}([5]).
\]
\end{exercise}

\begin{exercise}[1]
Find the number of maximal chains in

\[
\mathcal{P}([5]).
\]
\end{exercise}

\begin{exercise}[2]
Use Mantel's theorem to determine the largest possible number of edges
in a triangle-free graph on \(11\) vertices.
\end{exercise}

\begin{exercise}[2]
Construct an extremal triangle-free graph on \(11\) vertices.
\end{exercise}

\begin{exercise}[2]
Compute the number of edges in

\[
T_3(8).
\]
\end{exercise}

\begin{exercise}[2]
Use Tur\'an's theorem to determine

\[
\operatorname{ex}(8,K_4).
\]
\end{exercise}

\begin{exercise}[2]
Let

\[
\mathcal{F}
\subseteq
\binom{[9]}{3}
\]

be intersecting. Use the Erd\H{o}s--Ko--Rado theorem to bound
\(|\mathcal{F}|\).
\end{exercise}

\begin{exercise}[2]
Construct an intersecting family attaining the bound in the previous
exercise.
\end{exercise}

\begin{exercise}[2]
Use Sperner's theorem to determine the largest antichain in

\[
\mathcal{P}([8]).
\]
\end{exercise}

\begin{exercise}[2]
Explain why every level

\[
\binom{[n]}{k}
\]

of the Boolean lattice is an antichain.
\end{exercise}

\begin{exercise}[2]
Prove that the Boolean lattice has exactly

\[
n!
\]

maximal chains.
\end{exercise}

\begin{exercise}[2]
Compute the Lubell function of the middle level

\[
\binom{[n]}{\lfloor n/2\rfloor}.
\]
\end{exercise}

\begin{exercise}[3]
Prove Mantel's theorem using the degree argument given in this section.
\end{exercise}

\begin{exercise}[3]
Give a second proof of Mantel's theorem using a different method.
\end{exercise}

\begin{exercise}[3]
Prove that the balanced complete bipartite graph maximizes the product

\[
ab
\]

subject to

\[
a+b=n.
\]
\end{exercise}

\begin{exercise}[3]
Show that Mantel's theorem is the \(r=2\) case of Tur\'an's theorem.
\end{exercise}

\begin{exercise}[3]
Determine the exact number of edges in

\[
T_r(n)
\]

in terms of the quotient and remainder when \(n\) is divided by \(r\).
\end{exercise}

\begin{exercise}[3]
Prove Sperner's theorem using the LYM inequality.
\end{exercise}

\begin{exercise}[3]
Prove the LYM inequality by double counting maximal chains.
\end{exercise}

\begin{exercise}[3]
Explain why a full star in

\[
\binom{[n]}{k}
\]

is intersecting.
\end{exercise}

\begin{exercise}[3]
Compute the size of a full star in

\[
\binom{[n]}{k}.
\]
\end{exercise}

\begin{exercise}[3]
Using the random-subset argument, optimize

\[
pn-p^2m
\]

over \(p\).
\end{exercise}

\begin{exercise}[3]
Apply the resulting bound to a graph of average degree \(d\).
\end{exercise}

\begin{exercise}[3]
Give an example of an extremal problem where the extremal configuration
is not unique.
\end{exercise}

\begin{exercise}[3]
Explain the difference between exact extremal results and stability
results.
\end{exercise}

\begin{exercise}[4]
Study a proof of Tur\'an's theorem.
\end{exercise}

\begin{exercise}[4]
Investigate the Erd\H{o}s--Stone theorem and explain why chromatic
number determines asymptotic extremal density.
\end{exercise}

\begin{exercise}[4]
Study the Erd\H{o}s--Ko--Rado theorem and one proof using cyclic
permutations or shifting.
\end{exercise}

\begin{exercise}[4]
Investigate \(t\)-intersecting families and the Ahlswede--Khachatrian
complete intersection theorem.
\end{exercise}

\begin{exercise}[4]
Study compression and shifting methods in extremal set theory.
\end{exercise}

\begin{exercise}[4]
Investigate the polynomial method in a finite-set intersection problem.
\end{exercise}

\begin{exercise}[4]
Study a supersaturation theorem for triangles.
\end{exercise}

\begin{exercise}[4]
Investigate the triangle removal lemma.
\end{exercise}

\begin{exercise}[4]
Study an extremal problem for \(3\)-uniform hypergraphs.
\end{exercise}

\begin{exercise}[4]
Investigate a probabilistic proof of the existence of graphs with high
girth and high chromatic number.
\end{exercise}

%%%%%%%%%%%%%%%%%%%%%%%%%%%%%%%%%%%%%%%%%%%%%%%%%%%%%%%%%%%%
\subsection{Conceptual Summary}
\label{subsec:extremal-combinatorics-summary}
%%%%%%%%%%%%%%%%%%%%%%%%%%%%%%%%%%%%%%%%%%%%%%%%%%%%%%%%%%%%

Extremal combinatorics asks for the largest or smallest discrete
structures satisfying restrictions.

In graph theory,

\[
\boxed{
\operatorname{ex}(n,H)
}
\]

denotes the maximum number of edges in an \(n\)-vertex graph containing
no copy of \(H\).

Mantel's theorem gives

\[
\boxed{
\operatorname{ex}(n,K_3)
=
\left\lfloor\frac{n^2}{4}\right\rfloor.
}
\]

Tur\'an's theorem generalizes this to complete graphs:

\[
\boxed{
\operatorname{ex}(n,K_{r+1})
=
|E(T_r(n))|.
}
\]

The Erd\H{o}s--Stone theorem determines the leading asymptotic density
for general non-bipartite forbidden graphs.

In extremal set theory, the Erd\H{o}s--Ko--Rado theorem gives

\[
\boxed{
|\mathcal{F}|
\leq
\binom{n-1}{k-1}
}
\]

for intersecting \(k\)-uniform families when \(n\geq2k\).

Sperner's theorem gives

\[
\boxed{
|\mathcal{F}|
\leq
\binom{n}{\lfloor n/2\rfloor}
}
\]

for antichains in the Boolean lattice.

The LYM inequality strengthens this through

\[
\boxed{
\sum_{A\in\mathcal{F}}
\frac{1}{\binom{n}{|A|}}
\leq1.
}
\]

Important extremal methods include

\[
\boxed{
\text{double counting},
\quad
\text{averaging},
\quad
\text{compression},
\quad
\text{linear algebra},
\quad
\text{probability}.
}
\]

The probabilistic method uses

\[
\boxed{
\mathbb{P}(\text{desired property})>0
}
\]

to prove existence.

Stability theory studies structures that are close to extremal.

Thus extremal combinatorics seeks not only optimal numerical bounds but
also the structural reasons those bounds occur.

%%%%%%%%%%%%%%%%%%%%%%%%%%%%%%%%%%%%%%%%%%%%%%%%%%%%%%%%%%%%
\subsection{Section Connections}
\label{subsec:extremal-combinatorics-connections}
%%%%%%%%%%%%%%%%%%%%%%%%%%%%%%%%%%%%%%%%%%%%%%%%%%%%%%%%%%%%

\chapterconnection{
Extremal Combinatorics builds on Basic Counting and Enumerative
Combinatorics but changes the central question from counting every
admissible structure to determining the largest or smallest structure
that can satisfy a restriction. The next section, Algebraic
Combinatorics, introduces algebraic methods for studying discrete
objects and their symmetries. Probabilistic Combinatorics develops the
probabilistic method and random structures more systematically. Graph
Theory provides many of the central forbidden-subgraph problems.
Ramsey Theory studies the point at which avoidance becomes impossible.
Design Theory uses extremal bounds on highly regular incidence systems,
and Matroid Theory extends notions of independence and extremality to a
broad abstract setting.
}

%%%%%%%%%%%%%%%%%%%%%%%%%%%%%%%%%%%%%%%%%%%%%%%%%%%%%%%%%%%%
% SECTION SOURCE TRACKING
%%%%%%%%%%%%%%%%%%%%%%%%%%%%%%%%%%%%%%%%%%%%%%%%%%%%%%%%%%%%

%------------------------------------------------------------
% 6.6 Extremal Combinatorics
%------------------------------------------------------------
%
% Source basis:
%   - Standard extremal graph theory.
%   - Standard extremal set theory.
%   - Standard probabilistic and compression methods.
%
% Used for:
%   - extremal functions;
%   - forbidden configurations;
%   - triangle-free graphs;
%   - Mantel's theorem;
%   - balanced complete bipartite graphs;
%   - Turan graphs;
%   - Turan's theorem;
%   - Erdős--Stone theorem;
%   - bipartite extremal problems;
%   - intersecting families;
%   - Erdős--Ko--Rado theorem;
%   - stars in set systems;
%   - Boolean lattice;
%   - antichains;
%   - Sperner's theorem;
%   - maximal chains;
%   - Lubell function;
%   - LYM inequality;
%   - Dilworth's theorem;
%   - Mirsky's theorem;
%   - compression and shifting;
%   - averaging;
%   - average degree;
%   - double counting;
%   - probabilistic method;
%   - expectation method;
%   - alteration method;
%   - random-graph existence arguments;
%   - hypergraph extremal problems;
%   - t-intersecting families;
%   - forbidden intersection sizes;
%   - linear algebra method;
%   - polynomial method;
%   - stability;
%   - supersaturation;
%   - removal lemmas;
%   - exact versus asymptotic extremal results;
%   - extremal configurations;
%   - weighted extremal problems;
%   - applications and exercises.
%
% Notes:
%   - Detailed random-structure methods are developed in
%     Probabilistic Combinatorics.
%   - Graph-theoretic fundamentals are developed in the dedicated Graph
%     Theory section.
%   - Ramsey Theory develops unavoidable configurations.
%   - Hypergraph extremal theory is introduced only as a preview here.
%
%%%%%%%%%%%%%%%%%%%%%%%%%%%%%%%%%%%%%%%%%%%%%%%%%%%%%%%%%%%%